\documentclass[Total.tex]{subfiles}
\usepackage{booktabs}
\usepackage{array}
\begin{document}
	\begin{prob}
		(Das Problem des Handluingsereisenden/Traveling Salesman Problem) Given $n$ cities, with the distance between them. Then, we want to find the shortest cycle, which goes through each cities.
		
		Solution 1: With a start point, then at most $\frac{(n-1)!}{2}$ possibilities.
	\end{prob}
	\begin{prob}
		(Calculation of Pressure of Network) The calculation of graph of networks
	\end{prob}
	\chapter{Fundamentals}
		\subsection{Programs}
			The following list:
			\begin{table}[h]
				\centering
				\caption{Struktur und Bestandteile der Programmiersprache}
				\vspace{2mm}
				\begin{tabular}{@{}ll@{}}
					\toprule
					\textbf{Einheit} & \textbf{Bestandteile} \\ \midrule
					\textbf{Block} & \{ Programm \} \\ \addlinespace
					\textbf{Funktion} & Funktionskopf Block \\ \addlinespace
					\textbf{Funktionskopf} & Rückgabetyp Funktionsname ( Argumente ) \\ \addlinespace
					\textbf{Anweisung} & 
					\begin{tabular}[c]{@{}l@{}}
						; \\
						Ausdruck ; \\
						if-Bedingung \\
						while-Schleife \\
						for-Schleife \\
						Block
					\end{tabular} \\ \addlinespace
					\textbf{Ausdruck} & 
					\begin{tabular}[c]{@{}l@{}}
						Konstante \\
						Variablenauswertung \\
						Funktionsaufruf \\
						Operatoranwendung ( Ausdrücke )
					\end{tabular} \\ \addlinespace
					\textbf{Programm} & Variablendeklarationen, Funktionen, Anweisungen \\ \bottomrule
				\end{tabular}
			\end{table}
		\section{Logic}
			Omitted.
			\begin{gather*}
				NAND(a,b):=\neg(a\wedge b)=\neg a\vee \neg b\\
				Nor(a,b):=\neg(a\vee b)=\neg a\wedge \neg b\\
				Xor(a,b):=a\oplus b:=a\not=b
			\end{gather*}
			With C-codes:
			\begin{table}[h]
				\centering
				\caption{Logische Operatoren und Syntaxvergleich}
				\vspace{2mm}
				\begin{tabular}{@{}llll@{}}
					\toprule
					\textbf{Logik} & \textbf{Schriftform} & \textbf{C-Syntax} & \textbf{Assoziativität} \\ \midrule
					Negation       & $\neg$               & \texttt{!}        & von rechts nach links   \\
					Gleich         & $=$                  & \texttt{==}       & von links nach rechts   \\
					Ungleich       & $\neq$               & \texttt{!=}       & von links nach rechts   \\
					Und            & $\wedge$             & \texttt{\&\&}     & von links nach rechts   \\
					Oder           & $\vee$               & \texttt{||}       & von links nach rechts   \\
					Zuweisung      & $\leftarrow$         & \texttt{=}        & von rechts nach links   \\ \bottomrule
				\end{tabular}
			\end{table}
			\begin{Def}
				内容...
			\end{Def}
			
			\begin{Def}
				The output is logically false, iff it is $0$.
				
				The definition is well-defined for the output with image except $0$ and $1$.
			\end{Def}
			
			\begin{Def}
				A \textbf{program} consists of a sequence of commands.
				
				The most simple command is $;$.
				
				And 
			\end{Def}
			\subsection{If-Else}
		\section{Set}
			Omitted
		\section{Bit Operations}
			\begin{Def}
				A \textbf{bit} is the basic unit of information. 
			\end{Def}
			And:
			\begin{Def}
				Define the 
			\end{Def}
			Especially:
			\begin{Def}
				A word is a sequence of bits:With length $n$ called $n$-bit.
				\begin{itemize}
					\item A \textbf{byte} is a $8$-bit;
					\item A \textbf{nibble} is a $4$-bit;
				\end{itemize}
			\end{Def}
			\subsection{Logic Operations(for words)}
			\begin{gather*}
				\sim w:=(\neg wA_{n-1},\dots,\neg w_0)\\
				v\&w:=(v_{n-1}\wedge w_{n-1},\dots,v_0\wedge w_0)\\
				v|w:=(v_{n-1}\vee w_{n-1},\dots,v_0\vee w_0)\\
				v\hat{ }w:=(v_{n-1}\oplus w_{n-1}\dots,v_0\oplus w_0)		
			\end{gather*}
			\textbf{Remark} In C-program the "int" is not universal: (For example)It depends on $\#include<stdint.h>$.
			\begin{Def}
				(Logic Shifting) Define te operation:
				\begin{gather*}
					w<<m:=(w_{n-1-m},\dots,w_0,0,\dots,0)\\
					w>>m:=(0,\dots,0,w_{n-1},\dots,w_m)
				\end{gather*}
				Especially, $m=0,w<<(-)=id,w>>(-)=id$.
			\end{Def}
			\begin{Def}
				($0$-Word) A $0$-word is defined as the word with all bits are $0$.
			\end{Def}
			\subsection{Functions}
			\begin{table}[h]
				\centering
				\caption{Bitweise Operatoren und Zuweisungen in C}
				\vspace{2mm}
				\begin{tabular}{@{}lll@{}}
					\toprule
					\textbf{Operator} & \textbf{C-Syntax} & \textbf{Assoziativität} \\ \midrule
					Bitweise Negation & \texttt{\textasciitilde} & von rechts nach links \\
					Bitweiser Shift   & \texttt{<<, >>} & von links nach rechts \\
					Bitweises Und     & \texttt{\&} & von links nach rechts \\
					Bitweises XOR     & \texttt{\textasciicircum} & von links nach rechts \\
					Bitweises Oder    & \texttt{|} & von links nach rechts \\ \addlinespace
					Zuweisung         & \begin{tabular}[c]{@{}l@{}} \texttt{=, \&=, \textasciicircum=, |=,} \\ \texttt{<<=, >>=} \end{tabular} & von rechts nach links \\ \bottomrule
				\end{tabular}
			\end{table}
			Then define:
			\begin{Def}
				\begin{itemize}
					\item Set $i$ Bit to $0$:
					\begin{gather*}
						w\&=\sim(1<<i)
					\end{gather*}
					\item Set $i$ Bit to $1$:
					\begin{gather*}
						w|=1<<i
					\end{gather*}
					\item Bit Flip of $i$:
					\begin{gather*}
						w^{ }=1<<i
					\end{gather*}
					\item Testing $i$ bit is $0$ or $1$
					\begin{gather*}
						w \& (1 << i)\qquad \mbox{equivalent to}\qquad (w>>i)\&1
					\end{gather*}
				\end{itemize}
			\end{Def}
		
		\section{Natural Numbers}
			\subsection{Construction of Natural Number}
				We construct the natural number through bytes: Define the function
				\begin{gather*}
					s_n(w):=\begin{cases}
						(01)&w=(w_0)=(0),n=1\\
						(10)&w=(w_0)=(1),n=1\\
						(w_{n-1}\wedge w_{n-1}',w_{n-1}\oplus w_{n-1}',w_{n-2}',\dots,w_0')&w'=s_{n-1}(w_{n-2},\dots,w_0)
					\end{cases}
				\end{gather*}
				This actually defines
				\begin{gather*}
					s(-):\cup \{Seq_n()\}
				\end{gather*}
				which is actuarlly $s(n)=n+1$.
				\begin{proposition}
					There exists no $n$-bit word $w$ with $s_n(w)=0$.
				\end{proposition}
				\begin{proof}
					内容...
				\end{proof}
				This implies, that $0$ is the initial object.
				
				\begin{proposition}
					$s_n(w):=(1,q_{n-1},\dots,q_0)\Rightarrow q_i=0$
				\end{proposition}
				\begin{proof}
					内容...
				\end{proof}
				\begin{proposition}
					The followings are equivalent:
					\begin{gather*}
						w=(1,\dots,1)\Leftrightarrow s_n(w)=(1,0,\dots,0)
					\end{gather*}
				\end{proposition}
				\begin{proof}
					内容...
				\end{proof}
				This gives a correspondence:
				\begin{gather*}
					\{f\in 2^\NN|f\mbox{ is finite}\}\Leftrightarrow \NN
				\end{gather*}
				\begin{proposition}
					(Injectivity of Function $s_n$) $s_n$ is injective. 
				\end{proposition}
				\begin{proof}
					内容...
				\end{proof}
			\subsection{Order Relation}
				\begin{Def}
					Suppose $v$ be $m$-bit word and $w$ be $n$-bit word.
				\end{Def}
				
				\begin{corollary}
					Two words are equivalent, 
					\begin{gather*}
						v\sim w\Leftrightarrow [v]=[w]
					\end{gather*}
				\end{corollary}
				\begin{proof}
					内容...
				\end{proof}
				
				\begin{Def}
					内容...
				\end{Def}
				\begin{proposition}
					内容...
				\end{proposition}
				\begin{proof}
					内容...
				\end{proof}
				
				\begin{proposition}
					内容...
				\end{proposition}
				\begin{proof}
					内容...
				\end{proof}
				
				\begin{Def}
					内容...
				\end{Def}
		\section{Integers of Fixed Byte}
			\subsection{$b$-adic representation}
			\begin{Def}
				Define the $b$-adic representation of fixed byte:
				\begin{gather*}
					B(w)=\sum_{i=0}^{n-1} w_i b_i
				\end{gather*}
				Mostly we have $b_i=b^i$, so $B(w)$ is a 
			\end{Def}
			Then, recursively:
			
			\begin{proposition}
				(Horner-Schema) 
				\begin{gather*}
					\sum w_i b_i=b\dots(b(bw_{n-1}+w_{n-2})+w_{n-3}\dots)+w_0
				\end{gather*}
			\end{proposition}
			\begin{proof}
				Omitted.
			\end{proof}
			
			\begin{example}
				$10$-adic: 
				
				$16$-adic:
				
				$2$-adic:
				
				Define
				\begin{gather*}
					U(w)=\sum_{i=0}^{n-1} w_i 2^i\qquad \mathbb{U}=\NN\cap [0,2^n)
				\end{gather*}
				Then,
				\begin{gather*}
					U(w)=\sum_{i=0}^{n-1} w_i 2^i=\sum_{j=0}^{m-1}\sum_{i=0}^3 w_{4j+i}2^{4j+i}=\\
					\qquad\qquad =\sum_{j=0}^{m-1} 2^{4j}\sum_{i=0}^3 w_{4j+i}2^i\\
					\qquad\qquad =\sum_{j=0}^{m-1} z_j 16^h\qquad z_j=U(w_{4j+3},\dots,w_{4j})
				\end{gather*}
			\end{example}
			
			\subsection{Additions}
			\begin{Def}
				With adic representation as above, we have
				\begin{gather*}
					u_n(w)=\\
					v_n(v,w)=
				\end{gather*}
			\end{Def}
			\begin{proposition}
				
			\end{proposition}
			\begin{proof}
				内容...
			\end{proof}
			
			\begin{Def}
				Suppose $v,w$ be $n$-bit words, then define the \textbf{Two's complement}
				\begin{gather*}
					-_uw:=\sim w+_u 1
				\end{gather*}
				Then, define the operator
				\begin{gather*}
					v-_u w:=v+_u (-_u w)
				\end{gather*}
			\end{Def}
			\begin{proposition}
				内容...
			\end{proposition}
			\begin{proof}
				内容...
			\end{proof}
			
			\begin{corollary}
				内容...
			\end{corollary}
			\begin{proof}
				内容...
			\end{proof}
			\subsection{Modifications}
			\subsection{Types in $C$}
		\section{The Natural Numbers}
			Omitted.
		\section{Complexivity}
		
			\subsection{Landau-Symbols}
			\begin{Def}
				(Landau-Symbol)
				\begin{gather*}
					f(x)=\mathcal{O}(g(x))\Leftrightarrow \exists c\exists x_0(c\in \RR,x_0\in \RR,\forall x(x\geq x_0\rightarrow f(x)\leq cg(x))\\
					f(x)=o(g(x))\Leftrightarrow \forall c(c>0\rightarrow \exists x_0(x_0\in \RR\wedge \forall x(x\geq x_0\rightarrow f(x)\leq cg(x))))\\
					f(x)=\Theta(g(x))\Leftrightarrow 
				\end{gather*}
			\end{Def}
			Some of elementary results:
			\subsection{C of Algorithms}
		\section{Error Analysis}
	\chapter{Solution of Linear System of Equations}
		\section{Matrices, Vectors and their Implementation}
		\section{Gauss Algorithm, and LR-Decomposition}
			\subsection{Gauss Algorithms}
				\begin{proposition}
					Consider $n\times n$ matrix. Ten, Gauss-algorithm has the time complexivity of
					\begin{gather*}
						\mathcal{O}(n^3)
					\end{gather*}
				\end{proposition}
				\begin{proof}
					内容...
				\end{proof}
			\subsection{LR-Decomposition}
				Given a $A\in M_n(F)$. During the Gauss-elimination , define:
				\begin{gather*}
					l_{ik}:=\frac{a_{ik}^{(k-1)}}{a_k^{(k-1)}}
				\end{gather*}
				Define
				\begin{gather*}
					[G_k]_{ij}:=\begin{cases}
						l_{ik}&j=k,ik+1,\dots,n\\
						0&sonst
					\end{cases}
				\end{gather*}
				Then, we have
				\begin{gather*}
					A^{(k)}=(I-G_k)\cdot A^{(k-1)}\\
					b^{(k)}=(I-G_k)\cdot b^{k-1}
				\end{gather*}
				By
				\begin{gather*}
					[G_k A^{(k-1)}]_{ij}=\sum_{m=1}^n [G_k]_{im}[A^{(k-1)}]_{mj}=[G_k]_{ik}[A^{k-1}]_{kj}
				\end{gather*}
				\begin{proposition}
					Suppose $A\in M_n(F)$, and Gauss-Elimination can be applied without row exchange. 
					\begin{itemize}
						\item $(I-G_{n-1})\cdot \dots \cdot (I-G_1)A=R$;
						\item $G_k\cdot G_l=0$, $k\leq l$;
						\item $(I-G_k)^{-1}=(I+G_k)$;
						\item $((I-G_{n-1})\cdot \dots \cdot (I-G_1))^{-1}=I+\sum G_k=L$.
					\end{itemize}
				\end{proposition}
				\begin{proof}
					内容...
				\end{proof}
				
				\begin{corollary}
					\begin{gather*}
						A=L\cdot R
					\end{gather*}
				\end{corollary}
				\begin{proof}
					内容...
				\end{proof}
				
				\begin{corollary}
					\begin{gather*}
						PA=LR
					\end{gather*}
				\end{corollary}
				\begin{proof}
					内容...
				\end{proof}
			\subsection{Pivotelement}
			
		\section{Cholesky-Decomposition}
			Consider the positive definite matrix.
			
			\begin{proposition}
				\begin{itemize}
					\item 
					\item 
				\end{itemize}
			\end{proposition}
			\begin{proof}
				内容...
			\end{proof}
			\begin{proposition}
				（Orthonormal Decomposition for s.p.d) 
				\begin{gather*}
					A=LL^T
				\end{gather*}
			\end{proposition}
			\begin{proof}
				\begin{itemize}
					\item 
					\item $\Rightarrow$ Consider the s.p.d matrix
					\begin{gather*}
						\begin{pmatrix}
							A_{n-1}&b\\
							b^T&a_{nn}
						\end{pmatrix}\qquad b\in \RR^{n-1},A_{n-1}\in M_{n-1}(\RR)
					\end{gather*}
					By induction we have $A_{n-1}=L_{n-1}L_{n-1}^T$. Then, set
					\begin{gather*}
						L=\begin{pmatrix}
							L_{n-1}&0\\
							c^T&\alpha
						\end{pmatrix}
					\end{gather*}
					We have
					\begin{gather*}
						LL^T=\begin{pmatrix}
							L_{n-1}&0\\
							c^T&\alpha
						\end{pmatrix}\begin{pmatrix}
							L_{n-1}^T&c\\
							0&\alpha
						\end{pmatrix}=\begin{pmatrix}
							A_{n-1}&L_{n-1}c\\
							c^TL_{n-1}^T&c^Tc+\alpha^2
						\end{pmatrix}
					\end{gather*}
					We just need to get the value:
					\begin{gather*}
						\begin{cases}
							L_{n-1}c=b\\
							\alpha=\sqrt{a_{nn}-c^Tc}
						\end{cases}
					\end{gather*}
					And $L_{n-1}$ is inversable.
				\end{itemize}
			\end{proof}
			And we could apply the following correspondence by
			\begin{gather*}
				A=LL^T
			\end{gather*}
			Finally, we have
			\begin{gather*}
				\sum_{j=1}^n (n-j+1)=\mathcal{O}(n^2)+\frac{1}{6}n^3
			\end{gather}
	\chapter{Sorting Algorthims}
		\section{Definitions}
		\begin{prob}
			Suppose there exists $n\in \NN$ numbers, $z_i\in \RR$. Find a permutation $\pi_1,\dots,\pi_n$, with
			\begin{gather*}
				内容...
			\end{gather*}
		\end{prob}
		\section{Bumble Sort}
			Firstly, $<$ is well-defined. Then, by the applying the following steps:
			\begin{gather*}
				w_i>w_{i+1}\rightarrow \mbox{swap},i++\\
				Elsewise\rightarrow i++
			\end{gather*}
			Until the end. And consider the group of $w_{n-1},\dots,w_1$. Finaally, we get the sequence with
			\begin{gather*}
				w_{n-1}<\dots<w_0
			\end{gather*}
			\subsection{Time Complexity}
				\begin{itemize}
					\item Worst-Case: $\mathcal{O}(n^2)$;
					\item Average-Case: $\mathcal{O}(n^2)$;
					\item Best Basee: $\mathcal{O}(n)$, if we have the extra process
					\begin{gather*}
						\mbox{no exchange}\rightarrow stop
					\end{gather*}
				\end{itemize}
		\section{Mergesort}
			The main method:
			\begin{itemize}
				\item Divide the set $S$ into two part $S_1,S_2$;
				\item Apply the process recursively for $S_1,S_3$;
				\item Merge;
			\end{itemize}
			\begin{lemma}
				Given $S_x=\{x_1<\dot<x_m\},S_y=\{y_1<\dots<y_n\}$. Then $S=S_x\cup S_y$ can be sorted with at most $m+n-1$ steps.
			\end{lemma}
			\begin{proof}
				By creating a list 
			\end{proof}
			\subsection{Time Complexity}
			\begin{proposition}
				The mergesort has the time complexity of  $\mathcal{O}(nlogn)$
			\end{proposition}
			\begin{proof}
				Suppose $n=2^m$, We have:
				\begin{gather*}
					A(n)=(n-1)+2A(\frac{n}{2})
				\end{gather*}
			\end{proof}
		\section{Quicksort}
			By applying following process:
			\begin{itemize}
				\item Select Pivot:
				\item Recursitve Call:
				\item End:
			\end{itemize}
			\subsection{Time complexity}
			
		\section{Heapsort}
		\subsection{Binanry Tree,Arborescence,Heap}	
		\begin{Def}
			内容...
		\end{Def}
	\chapter{Graph Algorithms}
		Consider: Following datas
		\begin{itemize}
			\item $n$ vertices $V=\{v_0,\dots,v_{n-1}\}$;
			\item $m$ edges $E=\{e_0,\dots,e_{m-1}\}$;
		\end{itemize}
		where 
		\begin{gather*}
			e_j=\{v_{x_j},v_{y_j}\}\qquad e_j=(v_{x_j,v_{y_j}})\qquad x_j,y_j\in \{0,\dots,n-1\}
		\end{gather*}
		And define:
		\begin{gather*}
			C_0:=0\\
			C_{i+1}:=C_i+|\delta(v_i)|;i=0,\dots,n-=1
		\end{gather*}
		Then obviously, we have $C_n=2m$.
		
		For directed graph, define
		\begin{gather*}
			C_0^+:=0\qquad \\
			C_{i+1}^+:=\qquad 
		\end{gather*}
		\begin{Def}
			The 
		\end{Def}
	\chapter{Structure of Data}
	
		\section{Elementary Data Structures}
		
			\subsection{Stacks and Queues}
			
			\subsection{Linked Lists}
			
			\subsection{Implementing pointers and objects}
			
			\subsection{Representing Root Trees}
			
		\section{Hash-Table}
		
		\section{Binary Search Trees}
		
		\section{Red-Black Trees}
		
		\section{Augmenting Data Structures}
		
		\section{B-Trees}
		
		\section{Binomial Heaps}
		
		\section{Fibonacci Heaps}
		
		\section{Data Structures for Disjointed Sets}
\end{document}